% Part: history
% Chapter: biographies 
% Section: kurt-goedel

\documentclass[../../../include/open-logic-section]{subfiles}

\begin{document}

\olfileid{his}{bio}{god}

\olsection{Kurt Gödel}

\olphoto{goedel-kurt}{Kurt G{\"o}del}

Kurt 哥德尔（库尔特·哥德尔）于 1906 年 4 月 28 日出生在奥匈帝国的布吕恩（今捷克共和国的布尔诺）。由于生性好奇且聪慧，年幼的小库尔特（Kurtele）经常被家人称为“Der kleine Herr Warum”（小问号先生）。他从小学起学业就非常优异，唯有数学未曾拿到最高分。哥德尔 因健康欠佳经常缺课，并且免修体育课。他在童年时期被诊断出患有风湿热。尽管医学评估得出了相反的结论，但他一生都坚信这永久性地损伤了他的心脏。

哥德尔 于 1924 年开始在维也纳大学学习，并于 1929 年完成博士学业。他最初打算学习物理学，但部分由于哲学家 Rudolf Carnap（鲁道夫·卡尔纳普）的影响，他的兴趣很快转向了数学，尤其是逻辑学。他在 Hans Hahn（汉斯·哈恩）指导下完成的博士论文，证明了带等词的一阶谓词逻辑的完备性定理 \citep{Godel1929}。仅仅一年后，他取得了其最著名的成果——第一和第二不完全性定理（发表于 \citealt{Godel1931}）。在维也纳期间，哥德尔 深度参与了维也纳学圈，这是一个包括 Carnap 在内的具有科学精神的哲学家团体，Carnap 的工作尤其受到了 哥德尔 成果的影响。

1938 年，哥德尔 与 Adele Nimbursky（阿黛尔·宁布尔斯基）结婚。他的父母并不满意：她不仅比 哥德尔 年长六岁而且已经离异，还曾在夜总会担任舞者。然而，社会压力并未影响 哥德尔，他们的婚姻一直很幸福，直到 哥德尔 去世。

1938 年纳粹德国吞并奥地利后，哥德尔 和 Adele 移居美国，他在新泽西州的普林斯顿高等研究院任职。尽管性格内向且有些古怪，哥德尔 在普林斯顿的时光充满了合作与成果。他在集合论、哲学和物理学领域发表了文章。值得注意的是，他与同在高等研究院的同事 Albert Einstein（阿尔伯特·爱因斯坦）建立了特别深厚的友谊。

晚年，哥德尔 的心理健康状况恶化。1977 年妻子住院后，她便无法再为他做饭了。一生都饱受心理健康问题困扰的 哥德尔 彻底陷入了偏执。由于极度害怕被下毒，哥德尔 拒绝进食。他于 1978 年 1 月 14 日在普林斯顿死于饥饿。

\begin{reading}
关于 哥德尔 生平的完整传记，可参阅 \citet{Dawson1997}。更多传记文章以及论述 哥德尔 对逻辑学和哲学贡献的文章，可参阅 \citet{Wang1990}、\citet{Baaz2011}、\citet{Takeuti2003} 和 \citet{Sigmund2007}。

哥德尔 的博士论文可见德文原版 \citep{Godel1929}。不完全性定理的原版文献是 \citep{Godel1931}。哥德尔 所有已发表和未发表的著作，以及部分通信选集，均有英文版，收录于他的《文集》\citet{Godel1986,Godel1990}。

关于 哥德尔 不完全性定理的详细阐述，参见 \citet{Smith2013}。关于这些定理的非形式哲学讨论，可参见 Mark Linsenmayer 的播客 \citep{Linsenmayer2014}。
\end{reading}

\end{document}